\begin{definition}[Successive literal bond assignments on the endpoint route]%
    \label{def:peps_torus_swept_route_bond_assignments}
    \lean{TNLean.PEPS.sweptPhysicalBRouteDirection,
        TNLean.PEPS.torusVacatingPlaquetteMove,
        TNLean.PEPS.torusSweptPhysicalRouteOperators}
    \leanok
    \uses{def:peps_torus_swept_physical_endpoint_route,
        def:peps_torus_directed_bond_update,def:peps_torus_swept_string_assignment}
    Assume that the torus periods are at least eight and seven.
    Begin with the four-endpoint string assignment and take twelve steps
    through the following thirteen plaquette bases, translated by an
    arbitrary vertex:
    \[
        \begin{aligned}
             & (1,1)\to(2,1)\to(2,2)\to(3,2)\to(4,2)\to(4,1)\to(4,0)     \\
             & \qquad\to(3,0)\to(2,0)\to(1,0)\to(0,0)\to(0,-1)\to(0,-2).
        \end{aligned}
    \]
    At a source plaquette write $B,V,T,L$ for its bottom rightward,
    right upward, top rightward and left upward transports.
    A right, upward, left or downward step respectively replaces its
    shared transport by
    \begin{align}
        TL B^{-1},\qquad VB L^{-1},\qquad T^{-1}VB,\qquad V^{-1}TL.
    \end{align}
    All other ordered bonds, including the partner strings, are retained.
    Let $u_n$ denote the resulting recursively defined assignment after
    $n$ steps. These are auxiliary coefficient assignments for the flux
    motion of \cite[\S6.6, local lines 2340--2423]{Schuch2010PEPS}; no
    complete physical braid is asserted by their definition.
\end{definition}

\begin{theorem}[Derived vacancies and moving based fluxes]%
    \label{thm:peps_torus_swept_route_vacancies}
    \lean{TNLean.PEPS.regularWalkHolonomy_torusVacatingPlaquetteMove,
        TNLean.PEPS.regularWalkHolonomy_torusSweptPhysicalRoute_target_vacancy,
        TNLean.PEPS.regularWalkHolonomy_torusSweptPhysicalRoute_movingFlux}
    \leanok
    \uses{def:peps_torus_swept_route_bond_assignments,
        thm:peps_torus_directed_bond_update,thm:peps_torus_plaquette_holonomy}
    Each of the four elementary coefficient updates makes its source
    clockwise plaquette holonomy equal to the identity. For the specified
    twelve-step assignment, the target plaquette is vacant in the actual
    preceding assignment at every step.
    The moving based flux equals $h$ at stages zero and one,
    $ghg^{-1}$ at stages two through eight, and
    $hghg^{-1}h^{-1}$ at stages nine through twelve.
    Thus the retained partner string affects the based representative.
\end{theorem}

\begin{proof}\leanok
    Clockwise plaquette transport is $B^{-1}V^{-1}TL$.
    Substitution of each displayed update gives the first assertion.
    Starting from the literal initial string assignment, evaluate the
    selected transport recursively at each stage. The target holonomies
    reduce to the identity, and the source holonomies reduce to the three
    displayed words. Injectivity of the translated finite coordinate
    envelope justifies these evaluations even at periodic seams.
\end{proof}

\begin{definition}[Actual final endpoint and partner-comparison walks]%
    \label{def:peps_torus_swept_route_partner_walks}
    \lean{TNLean.PEPS.torusSweptPhysicalRouteFinalEndpoint,
        TNLean.PEPS.torusSweptPhysicalRouteAJointLoop,
        TNLean.PEPS.torusSweptPhysicalRouteReunionConnector,
        TNLean.PEPS.torusSweptPhysicalRouteBReunionLoop}
    \leanok
    \uses{def:peps_torus_swept_route_bond_assignments,
        def:peps_torus_swept_patch_endpoint_walks}
    Relative to the translation vertex, the final endpoint bases in the
    order $A_2,A_1,B_1,B_2$ are
    \[
        (-1,1),\quad(3,1),\quad(0,-2),\quad(1,-2).
    \]
    Retain the actual original connector from $A_2$ to $A_1$.
    Compare the final neighbouring B partners along the rightward edge
    from $(0,-2)$ to $(1,-2)$.
    In each joint loop, traverse the partner loop along the connector
    first and then the loop at the chosen base. Consequently its
    holonomy is the first based flux times the transported partner flux.
\end{definition}

\begin{theorem}[Actual reunion holonomies of the literal route]%
    \label{thm:peps_torus_swept_route_reunion_holonomy}
    \lean{TNLean.PEPS.regularWalkHolonomy_torusSweptPhysicalRoute_finalEndpoints,
        TNLean.PEPS.regularWalkHolonomy_torusSweptPhysicalRoute_finalConnectorA,
        TNLean.PEPS.regularWalkHolonomy_torusSweptPhysicalRoute_transportedAPartner,
        TNLean.PEPS.regularWalkHolonomy_torusSweptPhysicalRoute_AJointLoop,
        TNLean.PEPS.regularWalkHolonomy_torusSweptPhysicalRoute_reunionConnector,
        TNLean.PEPS.regularWalkHolonomy_torusSweptPhysicalRoute_BReunionLoop,
        TNLean.PEPS.regularWalkHolonomy_torusSweptPhysicalRoute_BReunionLoop_conjClass,
        TNLean.PEPS.regularWalkHolonomy_torusSweptPhysicalRoute_BReunionLoop_eq_one_iff}
    \leanok
    \uses{def:peps_torus_swept_route_partner_walks,
        thm:peps_torus_swept_route_vacancies,thm:peps_regular_connector_holonomy,
        thm:peps_conjugacy_class_invariance}
    The final based endpoint fluxes are
    \begin{align}
        (g^{-1},\ g,\ hghg^{-1}h^{-1},\ h^{-1}).
    \end{align}
    The actual A connector has transport $h^{-1}g^{-1}$; hence the A1 loop
    transported to A2 has flux $ghgh^{-1}g^{-1}$, and their joint loop has
    flux $hgh^{-1}g^{-1}$. The final B reunion connector has identity
    transport and its partner-first joint loop has flux
    \begin{align}
        h(ghg^{-1}h^{-1})h^{-1}.
    \end{align}
    Its conjugacy class is that of the commutator, and its joint flux is
    trivial precisely when $g$ and $h$ commute.
    These assertions concern the recursively derived coefficient
    assignment and actual comparison walks. Physical implementation of
    every step of the route remains a separate assertion.
\end{theorem}

\begin{proof}\leanok
    Evaluate the four actual plaquette loops and both connectors from
    the derived final transports. Along a connector $p$, partner-loop
    transport is $H(p)^{-1}H(\ell)H(p)$. Concatenation reverses the
    multiplication order, giving the stated joint products. Conjugation
    preserves the conjugacy class, and a group commutator is the identity
    exactly when its two arguments commute.
\end{proof}
